\documentclass[11pt]{amsart}
\usepackage[T1]{fontenc}
\usepackage{lmodern,microtype,amsmath,amssymb,amsthm,mathtools,enumitem}
\usepackage[margin=1in]{geometry}
\usepackage{booktabs}
\usepackage[hidelinks]{hyperref}
\newtheorem{mainthm}{Theorem}
\renewcommand{\themainthm}{\Alph{mainthm}}
\newtheorem{theorem}{Theorem}[section]
\newtheorem{proposition}[theorem]{Proposition}
\newtheorem{lemma}[theorem]{Lemma}
\newtheorem{corollary}[theorem]{Corollary}
\theoremstyle{definition}
\newtheorem{definition}[theorem]{Definition}
\newtheorem{question}[theorem]{Question}
\theoremstyle{remark}
\newtheorem{remark}[theorem]{Remark}
\newcommand{\bits}{\{-1,1\}}
\newcommand{\E}{\mathbb E}
\renewcommand{\P}{\mathbb P}
\newcommand{\R}{\mathbb R}
\newcommand{\HH}{\mathsf H}
\newcommand{\I}{\mathsf I}
\newcommand{\wh}[1]{\widehat{#1}}
\newcommand{\eps}{\varepsilon}
\DeclareMathOperator{\Stab}{Stab}
\DeclareMathOperator{\Ei}{E_1}
\setlist{itemsep=3pt,topsep=5pt}
\setlength{\emergencystretch}{1em}
\raggedbottom
\hypersetup{pdftitle={The Fourier entropy-influence inequality for symmetric Boolean functions},pdfauthor={Yitzchak Shmalo}}
\title[Fourier entropy of symmetric functions]{The Fourier entropy--influence inequality for symmetric Boolean functions}
\author{Yitzchak Shmalo}
\address{Einstein Institute of Mathematics, The Hebrew University of Jerusalem, Jerusalem, Israel}
\email{yitzchak.shmalo@gmail.com}
\date{9 November 2026}
\subjclass[2020]{42C10, 60C05, 68Q17, 06E30, 33C10}
\keywords{Fourier entropy--influence conjecture, symmetric Boolean function, noise stability, Krawtchouk polynomial, modified Bessel function}
\usepackage{xurl}
\begin{document}
\begin{abstract}
We prove that every symmetric Boolean function $f$ satisfies
$\HH[f]\le\bigl(\log_2(4+e)+2+1/\ln 2\bigr)\I[f]<6.1908\,\I[f]$, improving the constant $12.04$ of O'Donnell, Wright and Zhou. The Fourier entropy of a symmetric function splits into the entropy of its level weights and a within-level term. The first is controlled for every Boolean function by a spectral-level estimate. For the second we prove that the derivative of a symmetric function has noise stability at most $e^{-v}I_0(v)$ times its squared norm, where $v$ is the variance of the noise, and evaluate the resulting entropy integral through a Frullani identity and two inequalities for modified Bessel functions. The AND and OR functions show that the optimal constant for symmetric functions is at least $4$. We also prove that the entropy cost of restricting one coordinate of a symmetric function is asymptotically at most $4.637$ times its influence, so the restriction step that would give constant $4.8723$ can fail in only finitely many dimensions.
\end{abstract}
\maketitle

\section{Introduction}

For $f:\bits^n\to\bits$ write $f=\sum_{S\subseteq[n]}\wh f(S)\chi_S$ with $\chi_S(x)=\prod_{i\in S}x_i$. By Parseval the squared coefficients form a probability distribution. Its entropy and the total influence are
\[
 \HH[f]=\sum_{S}\wh f(S)^2\log_2\frac{1}{\wh f(S)^2},\qquad
 \I[f]=\sum_S|S|\,\wh f(S)^2=\sum_{i=1}^n\P_x\bigl[f(x)\ne f(x^{\oplus i})\bigr],
\]
where $x^{\oplus i}$ is $x$ with its $i$-th coordinate flipped. The Fourier entropy--influence conjecture of Friedgut and Kalai~\cite{FK96} asserts that $\HH[f]\le C\,\I[f]$ for an absolute constant $C$. It remains open. It is known for several classes of functions, among them symmetric functions~\cite{OWZ11,KMS12,CKLS16,ACKSW21,Shm26}.

A function is \emph{symmetric} if $f(x)$ depends only on the Hamming weight $|x|=\#\{i:x_i=-1\}$. O'Donnell, Wright and Zhou~\cite{OWZ11} proved the conjecture for symmetric functions with constant $12.04$. Our main result lowers this constant by about half.

\begin{mainthm}\label{thm:main}
Every symmetric Boolean function $f$ satisfies
\[
 \HH[f]\le C_{\rm sym}\,\I[f],\qquad
 C_{\rm sym}=\log_2(4+e)+2+\frac{1}{\ln2}=6.1907873580\ldots<6.1908.
\]
\end{mainthm}

The proof has two parts. For symmetric $f$ the coefficient $\wh f(S)$ depends only on $k=|S|$, and the entropy splits as $\HH=L+B$, where $L$ is the entropy of the level weights $W_k=\sum_{|S|=k}\wh f(S)^2$ and $B=\sum_kW_k\log_2\binom nk$ measures the spread inside the levels. The bound $L\le\log_2(4+e)\I$ holds for every Boolean function; it is Lemma~2.1 of~\cite{Shm26}, and we recall its short proof in Section~\ref{sec:level}. The new part is the within-level bound $B<(2+1/\ln2)\I$. It rests on the following noise-stability estimate, which is of independent interest. Write $T_\lambda$ for the noise operator, $\Stab_\lambda[g]=\langle g,T_\lambda g\rangle=\sum_S\lambda^{|S|}\wh g(S)^2$, and $I_0$ for the modified Bessel function of order zero.

\begin{mainthm}\label{thm:stab}
Let $f$ be a symmetric Boolean function on $n\ge2$ variables, $N=n-1$, and let $g=D_nf$ be its derivative in the last coordinate, a function on $\bits^N$. For $0\le\lambda\le1$,
\[
 \Stab_\lambda[g]\le e^{-v}I_0(v)\,\|g\|_2^2,\qquad v=\frac{N(1-\lambda^2)}{4}.
\]
\end{mainthm}

Since $e^{-v}I_0(v)\sim(2\pi v)^{-1/2}$, this improves the derivative estimate used in~\cite{OWZ11} by a factor that tends to $2$ when $\lambda\to1$ and $v\to\infty$ together. The proof uses three facts: the derivative of a symmetric function is an alternating sum of indicators of Hamming slices; an alternating sum of selected values of a unimodal sequence is at most its largest value; and every atom of a sum of independent Bernoulli variables is bounded by Fourier inversion. The within-level term is then bounded by an integral of the envelope $e^{-v}I_0(v)$ against a harmonic-number kernel, which we show is less than $\ln4$ for every $N$. A Frullani identity evaluates the main part exactly, and two inequalities for $e^{-x}I_0(x)$ and $e^{-x}I_1(x)$ control the rest.

The constant $4$ is a natural lower bound.

\begin{mainthm}\label{thm:lower}
For the OR function on $n$ variables (and for AND),
\[
 \frac{\HH[\mathrm{OR}_n]}{\I[\mathrm{OR}_n]}=4-\frac{4-2/\ln2}{n}+O(2^{-n}).
\]
Hence the optimal constant for symmetric functions lies in $[4,\,6.1908]$.
\end{mainthm}

That OR approaches $4$ is classical; we include the exact expression because AND and OR are the maximizers among all symmetric functions in every dimension we computed (Section~\ref{sec:computations}).

The inductive approach to the conjecture asks for a coordinate whose restriction costs little entropy. For a coordinate $i$ let $f_{i\to\pm1}$ be the two restrictions and
\[
 \Delta_i[f]=\HH[f]-\tfrac12\bigl(\HH[f_{i\to1}]+\HH[f_{i\to-1}]\bigr).
\]
If $\Delta_1[f]\le K\,\I_1[f]$ for every symmetric $f$, then $\HH\le K\I$ for every symmetric $f$ by induction on $n$ (Lemma~\ref{lem:induction}). The best $K$ is at least $\Delta_1/\I_1$ for $\mathrm{OR}_3$, which equals $12-\tfrac92\log_23=4.8677\ldots$, so this route cannot reach $4$; in all dimensions $n\le16$ it is the largest value. We prove that the restriction step is much cheaper in high dimension.

\begin{mainthm}\label{thm:restriction}
\[
 \limsup_{n\to\infty}\ \sup_{f}\ \frac{\Delta_1[f]}{\I_1[f]}\le C_*<4.637,
\]
where the supremum is over nonconstant symmetric $f$ on $n$ variables and
\[
 C_*=\frac{1}{\ln2}\Bigl[\frac{14}{5}+\frac{89}{50}\,e^{-10}I_0(10)+\frac{48}{23\sqrt{40\pi}}+\frac{e^{-z}}{z}\Bigr]=4.6364236659\ldots,\qquad z=40\Bigl(\frac{23}{48}\Bigr)^2.
\]
\end{mainthm}

Consequently the restriction step with constant $4.8723$ can fail for at most finitely many $n$. An explicit cutoff, together with a finite computation below it, would give $\HH\le4.8723\,\I$ for every symmetric $f$. We do not have such a cutoff (Question~\ref{q:cutoff}).

The proof of Theorem~\ref{thm:restriction} uses Theorem~\ref{thm:stab} in place of a log-Sobolev inequality, a passage to a limiting degree distribution, and one pointwise majorant of the restriction-entropy integrand that is verified by outward-rounded interval arithmetic. All other steps are analytic. The computations are in the repository listed under Data availability.

\subsection*{Notation}
Logarithms are natural unless written $\log_2$. Entropies are in bits. For a symmetric function we write $F(w)\in\bits$ for its value on weight $w$, $\wh f_k$ for the common coefficient at level $k$, and $W_k=\binom nk\wh f_k^{\,2}$. We put
\[
 \Phi(x)=e^{-x}I_0(x),\qquad G(x)=e^{-x}I_1(x)\qquad(x\ge0),
\]
so that $\Phi(x)=\frac1\pi\int_0^\pi e^{-x(1-\cos\theta)}\,d\theta$ and $\Phi'=G-\Phi$. Finally $\Ei(x)=\int_x^\infty e^{-u}u^{-1}\,du$ and $H_N=\sum_{i\le N}1/i$.

\section{The entropy of a symmetric function}\label{sec:split}

For symmetric $f$ every set of size $k$ carries the same coefficient, so
\begin{equation}\label{eq:split}
 \HH[f]=\sum_{k=0}^n\binom nk\wh f_k^{\,2}\log_2\frac{1}{\wh f_k^{\,2}}
 =\underbrace{\sum_kW_k\log_2\frac1{W_k}}_{L}+\underbrace{\sum_kW_k\log_2\binom nk}_{B}.
\end{equation}
The level weights satisfy $\sum_kW_k=1$ and $\sum_kkW_k=\I[f]$.

For $n\ge2$ let $N=n-1$ and $g=D_nf$, where $D_nf(x)=\tfrac12\bigl(f(x,1)-f(x,-1)\bigr)$ for $x\in\bits^N$. Then $\wh g(T)=\wh f(T\cup\{n\})$, so $g$ is symmetric with
\begin{equation}\label{eq:derivative}
 W_j[g]=\binom Nj\wh f_{j+1}^{\,2}=\frac{(j+1)W_{j+1}}{n},\qquad \|g\|_2^2=\sum_jW_j[g]=\frac{\I[f]}n.
\end{equation}
Its values are $g(x)=\tfrac12\bigl(F(w)-F(w+1)\bigr)$ with $w=|x|$. They are nonzero exactly on the weights $w_0<w_1<\dots$ where $F$ changes value, and there they alternate in sign. Up to a global sign,
\begin{equation}\label{eq:alternating}
 g(x)=\sum_{i}(-1)^i\,\mathbf 1[|x|=w_i].
\end{equation}
In particular $|g|=g^2$, so $\E|g|=\|g\|_2^2$.

\section{The level entropy}\label{sec:level}

\begin{lemma}[{\cite[Lemma~2.1]{Shm26}}]\label{lem:level}
Every Boolean function satisfies $L\le\log_2(4+e)\,\I[f]$.
\end{lemma}

\begin{proof}
We may assume $f$ is not constant. Let $\eps=1-W_0>0$ and $Q(k)=W_k/\eps$ for $k\ge1$. Put $a=\log_2(1+e/4)$, so that $(2^a-1)2^{-ak}$, $k\ge1$, is a probability law. Gibbs' inequality against this law gives $\HH(Q)\le a\,\E_Qk-\log_2(2^a-1)$. Since $L=h_2(\eps)+\eps\HH(Q)$, $\eps\E_Qk=\I$ and $h_2(\eps)\le\eps\log_2(1/\eps)+\eps\log_2e$, we get
\[
 L\le a\I+\eps\bigl[\log_2(1/\eps)+\log_2e-\log_2(e/4)\bigr]=a\I+\eps\bigl[\log_2(1/\eps)+2\bigr].
\]
Let $p=\min\{\P[f=1],\P[f=-1]\}$, so $\eps=4p(1-p)$. The edge-isoperimetric inequality gives $\I\ge2p\log_2(1/p)$, and
\[
 2p\log_2\frac1p-\eps\Bigl(1+\tfrac12\log_2\frac1\eps\Bigr)=2p\bigl[p\log_2\tfrac1p+(1-p)\log_2(1-p)\bigr]\ge0,
\]
because the bracket is concave on $[0,\tfrac12]$ (its second derivative is $(2p-1)/(p(1-p)\ln2)$) and vanishes at both ends. Hence $\eps[\log_2(1/\eps)+2]\le2\I$ and $L\le(a+2)\I=\log_2(4+e)\I$.
\end{proof}

\section{Noise stability of derivatives}\label{sec:stab}

\begin{lemma}\label{lem:unimodal}
Let $(a_y)_{y\ge0}$ be a nonnegative unimodal sequence and $y_0<y_1<\dots<y_m$ integers. Then $\bigl|\sum_i(-1)^ia_{y_i}\bigr|\le\max_ya_y$.
\end{lemma}

\begin{proof}
Write $a_y=\int_0^\infty\mathbf 1[a_y>t]\,dt$. For each $t$ the set $\{y:a_y>t\}$ is an interval, so the selected indices in it form a block of consecutive $i$, whose alternating sum is $0$ or $\pm1$. The integrand vanishes for $t\ge\max a$. Integrate.
\end{proof}

\begin{lemma}\label{lem:atom}
Let $Z$ be a sum of $N$ independent Bernoulli variables with parameters in $\{p,1-p\}$, and $v=Np(1-p)$. Then $Z$ has a unimodal law and $\P[Z=y]\le\Phi(v)$ for every integer $y$.
\end{lemma}

\begin{proof}
Each Bernoulli law is log-concave and convolution preserves log-concavity of sequences supported on intervals, so the law of $Z$ is log-concave and in particular unimodal. Each summand has characteristic function of modulus $\bigl(1-4p(1-p)\sin^2(\theta/2)\bigr)^{1/2}$. Using $1-u\le e^{-u}$,
\[
 |\E e^{i\theta Z}|\le\exp\bigl(-2v\sin^2(\theta/2)\bigr)=e^{-v(1-\cos\theta)}.
\]
Fourier inversion gives $\P[Z=y]\le\frac1{2\pi}\int_{-\pi}^\pi|\E e^{i\theta Z}|\,d\theta\le\Phi(v)$.
\end{proof}

\begin{proof}[Proof of Theorem~\ref{thm:stab}]
Let $p=(1-\lambda)/2$. Given $x$ with $|x|=r$, the weight of a $\lambda$-correlated copy $y$ of $x$ is $Z_r=\mathrm{Bin}(r,1-p)+\mathrm{Bin}(N-r,p)$, of variance $Np(1-p)=v$. By~\eqref{eq:alternating},
\[
 \Stab_\lambda[g]=\E_x\Bigl[g(x)\sum_i(-1)^i\,\P\bigl[Z_{|x|}=w_i\bigr]\Bigr].
\]
By Lemmas~\ref{lem:unimodal} and~\ref{lem:atom} the inner sum has absolute value at most $\Phi(v)$. Hence $\Stab_\lambda[g]\le\Phi(v)\,\E|g|=\Phi(v)\|g\|_2^2$.
\end{proof}

\begin{remark}
For even $N$ the atom bound before the step $1-u\le e^{-u}$ is attained at $r=y=N/2$. The constant $\sup_v\sqrt v\,\Phi(v)=0.46882\ldots$ is the sharp uniform concentration constant for Bernoulli sums of Baillon, Cominetti and Vaisman~\cite{BCV16}; we use the full function $\Phi$ rather than this constant.
\end{remark}

\section{The within-level entropy}\label{sec:within}

\begin{lemma}\label{lem:reduction}
Let $f$ be symmetric on $n\ge2$ variables with $\I[f]>0$, and $N=n-1$. Then
\[
 B\le\frac{1+C_N}{\ln2}\,\I[f],\qquad
 C_N=\int_0^1\frac{\Phi\bigl(Nt(2-t)/4\bigr)-(1-t)^N}{t}\,dt.
\]
\end{lemma}

\begin{proof}
Using $\binom nk\le(en/k)^k$,
\[
 B\ln2\le\sum_{k\ge1}kW_k\Bigl(1+\ln\frac nk\Bigr)=\I\Bigl(1+\E_q\ln\frac n{j+1}\Bigr),\qquad q_j=\frac{(j+1)W_{j+1}}{\I},\ 0\le j\le N.
\]
By~\eqref{eq:derivative}, $q_j=W_j[g]/\|g\|_2^2$. Comparing the sum with an integral, $\ln\frac{n}{j+1}\le H_N-H_j=\int_0^1\frac{\rho^j-\rho^N}{1-\rho}\,d\rho$. Therefore
\[
 \E_q\ln\frac n{j+1}\le\int_0^1\frac{s(\rho)-\rho^N}{1-\rho}\,d\rho,\qquad s(\rho)=\sum_jq_j\rho^j=\frac{\Stab_\rho[g]}{\|g\|_2^2}.
\]
By Theorem~\ref{thm:stab}, $s(\rho)\le\Phi\bigl(N(1-\rho^2)/4\bigr)$. Substituting $\rho=1-t$, so that $1-\rho^2=t(2-t)$, gives the claim. The integrand of $C_N$ is $O(1)$ near $t=0$.
\end{proof}

We need three facts about $\Phi$ and $G$.

\begin{lemma}\label{lem:bessel}
\begin{enumerate}[label=\textup{(\roman*)}]
\item $\sqrt w\,\Phi(w)>1/\sqrt{2\pi}$ for every $w\ge\tfrac12$.
\item $J(X):=\int_0^Xu\,\bigl(-\Phi'(u)\bigr)\,du=XG(X)$ for every $X\ge0$.
\item $G(X)<1/\sqrt{2\pi X}$ for every $X>0$.
\end{enumerate}
\end{lemma}

\begin{proof}
(i) Let $h(w)=\sqrt w\,\Phi(w)$. Substituting $I_0=e^w\Phi$ into the Bessel equation $I_0''+I_0'/w-I_0=0$ gives
\[
 h''+2h'+\frac{h}{4w^2}=0.
\]
At a stationary point $h''=-h/(4w^2)<0$, so $h$ has no local minimum in $(0,\infty)$. At the left end $h(\tfrac12)\ge e^{-1/2}/\sqrt2>1/\sqrt{2\pi}$, since $I_0\ge1$ and $\pi>e$. As $w\to\infty$, the substitution $\theta=y/\sqrt w$ and dominated convergence (with $1-\cos\theta\ge2\theta^2/\pi^2$ on $[0,\pi]$) give $h(w)\to\frac1\pi\int_0^\infty e^{-y^2/2}dy=1/\sqrt{2\pi}$. If $h(w_0)\le1/\sqrt{2\pi}$ for some $w_0>\tfrac12$, then $h$ restricted to some $[\tfrac12,W]$ with $W>w_0$ attains its minimum at an interior point: if $h(w_0)<1/\sqrt{2\pi}$ choose $W$ with $h(W)>h(w_0)$, and if equality holds take any $W>w_0$, since then $h\ge h(w_0)$ on $[\tfrac12,\infty)$ by the first case. This interior minimum is a local minimum, a contradiction.

(ii) From $I_0'=I_1$ and $I_1'=I_0-I_1/x$ we get $\Phi'=G-\Phi$ and $G'=\Phi-G-G/x$. Hence $(xG)'=x(\Phi-G)=-x\Phi'$, and $xG(x)\to0$ as $x\to0$.

(iii) Integrating $I_1(X)=\frac1\pi\int_0^\pi e^{X\cos\theta}\cos\theta\,d\theta$ by parts and substituting $u=1-\cos\theta$,
\[
 G(X)=\frac X\pi\int_0^\pi e^{-X(1-\cos\theta)}\sin^2\theta\,d\theta=\frac X\pi\int_0^2e^{-Xu}\sqrt{u(2-u)}\,du
 <\frac{\sqrt2X}{\pi}\int_0^\infty e^{-Xu}\sqrt u\,du=\frac1{\sqrt{2\pi X}}.\qedhere
\]
\end{proof}

\begin{proposition}\label{prop:CN}
$C_N<\ln4$ for every integer $N\ge1$. More precisely $C_1\le1$, $C_2\le59/64$, and for $N\ge3$
\[
 C_N<\ln4-\frac{2-\sqrt2}{\sqrt{\pi N}}+\frac1{2N}<\ln4.
\]
\end{proposition}

\begin{proof}
Split $C_N=T_1+T_2+T_3$ with
\[
 T_1=\int_0^1\frac{\Phi(\frac{Nt(2-t)}4)-\Phi(\frac{Nt}2)}{t}dt,\quad
 T_2=\int_0^1\frac{\Phi(\frac{Nt}2)-e^{-Nt}}{t}dt,\quad
 T_3=\int_0^1\frac{e^{-Nt}-(1-t)^N}{t}dt.
\]

\emph{The term $T_2$.} Write $\Phi(x)=\E_\theta e^{-x(1-\cos\theta)}$ with $\theta$ uniform on $[0,\pi]$. For each $\theta$, Frullani's integral gives $\int_0^\infty\bigl(e^{-Nt(1-\cos\theta)/2}-e^{-Nt}\bigr)\frac{dt}t=\ln\frac2{1-\cos\theta}$; the integrand is nonnegative, so by Tonelli
\[
 \int_0^\infty\frac{\Phi(Nt/2)-e^{-Nt}}{t}dt=\E_\theta\ln\frac{1}{\sin^2(\theta/2)}=-\frac4\pi\int_0^{\pi/2}\ln\sin u\,du=\ln4.
\]
On $t\ge1$ we have $Nt/2\ge\tfrac12$, so Lemma~\ref{lem:bessel}(i) gives $\int_1^\infty\Phi(Nt/2)\frac{dt}t\ge\int_1^\infty\frac{dt}{\sqrt{\pi N}\,t^{3/2}}=\frac2{\sqrt{\pi N}}$. Since $\int_1^\infty e^{-Nt}\frac{dt}t=\Ei(N)$,
\[
 T_2\le\ln4-\frac2{\sqrt{\pi N}}+\Ei(N).
\]

\emph{The term $T_1$.} The function $\Phi$ is completely monotone, so $-\Phi'$ is positive and decreasing. With $a=Nt(2-t)/4$ and $b=Nt/2$ we have $a\ge Nt/4$ and $b-a=Nt^2/4$, so $\Phi(a)-\Phi(b)\le\frac{Nt^2}4\bigl(-\Phi'(\frac{Nt}4)\bigr)$. By Lemma~\ref{lem:bessel}(ii) and~(iii),
\[
 T_1\le\int_0^1\frac{Nt}4\Bigl(-\Phi'\bigl(\tfrac{Nt}4\bigr)\Bigr)dt=\frac4NJ\bigl(\tfrac N4\bigr)=G\bigl(\tfrac N4\bigr)<\sqrt{\frac{2}{\pi N}}.
\]

\emph{The term $T_3$.} Since $\int_0^1\frac{1-(1-t)^N}{t}dt=H_N$ and $\int_0^1\frac{1-e^{-Nt}}tdt=\ln N+\gamma+\Ei(N)$,
\[
 T_3=H_N-\ln N-\gamma-\Ei(N)\ \ (\ge0).
\]

Adding, the $\Ei(N)$ terms cancel and $H_N-\ln N-\gamma<1/(2N)$ gives
\[
 C_N<\ln4+\frac{\sqrt2}{\sqrt{\pi N}}-\frac{2}{\sqrt{\pi N}}+\frac1{2N}.
\]
The last two terms together are negative when $N>\pi/\bigl(4(2-\sqrt2)^2\bigr)$, and this threshold is less than $3$ because $\pi<4<12(2-\sqrt2)^2$, the last inequality following from $\sqrt2<17/12$.

For $N=1$, $\Phi\le1$ gives $C_1\le\int_0^1\frac{1-(1-t)}tdt=1$. For $N=2$, $e^{-z}\le1-z+z^2/2$ together with $\E(1-\cos\theta)=1$ and $\E(1-\cos\theta)^2=\tfrac32$ gives $\Phi(v)\le1-v+\tfrac34v^2$. With $v=(1-\rho^2)/2$,
\[
 C_2\le\int_0^1\frac{\frac{1-\rho^2}2+\frac3{16}(1-\rho^2)^2}{1-\rho}\,d\rho=\frac34+\frac3{16}\cdot\frac{11}{12}=\frac{59}{64}.
\]
Both are below $\ln4$.
\end{proof}

\begin{proof}[Proof of Theorem~\ref{thm:main}]
For $n=1$ every Boolean function is constant or $\pm x_1$ and $\HH=0$. For $n\ge2$ and $\I>0$, Lemma~\ref{lem:reduction} and Proposition~\ref{prop:CN} give $B<\frac{1+\ln4}{\ln2}\I=(2+\frac1{\ln2})\I$, and Lemma~\ref{lem:level} gives $L\le\log_2(4+e)\I$. Add them using~\eqref{eq:split}.
\end{proof}

\begin{remark}
The two parts are not tight together. For majority on $n=2m+1$ variables, $B/\I\to2/\ln2$, so the within-level constant cannot go below $2.885$; and $L/\I\le2.52$ for every symmetric function with $n\le16$ (Section~\ref{sec:computations}), well below $\log_2(4+e)=2.748$. Numerically $C_N$ increases to $\ln4$: $C_1=0.6638$, $C_2=0.8911$, $C_3=0.9991$, $C_{10^6}=1.3855$.
\end{remark}

\section{AND and OR}\label{sec:lower}

\begin{proof}[Proof of Theorem~\ref{thm:lower}]
Let $f$ take the value $-1$ at exactly one point; this is OR or AND up to the signs of the inputs and the output, which change neither $\HH$ nor $\I$. Then $\wh f(\emptyset)=1-2^{1-n}$ and $\wh f(S)^2=4^{1-n}$ for $S\ne\emptyset$. Hence, with $W_0=(1-2^{1-n})^2$,
\[
 \HH[f]=W_0\log_2\frac1{W_0}+(2^n-1)\,4^{1-n}\cdot2(n-1),\qquad \I[f]=n\,2^{1-n}.
\]
The second term divided by $\I$ is $(2-2^{1-n})\frac{2(n-1)}n=4-\frac4n+O(2^{-n})$, and $W_0\log_2(1/W_0)=\frac{2^{2-n}}{\ln2}\bigl(1+O(2^{-n})\bigr)$ contributes $\frac{2}{n\ln2}+O(2^{-n})$.
\end{proof}

At $n=1024$ the ratio is $3.998911513751\ldots$.

\section{The restriction step}\label{sec:restriction}

\begin{lemma}\label{lem:induction}
Suppose $\Delta_1[f]\le K\,\I_1[f]$ for every symmetric $f$ on at most $m$ variables. Then $\HH[f]\le K\,\I[f]$ for every symmetric $f$ on at most $m$ variables.
\end{lemma}

\begin{proof}
Restrictions of a symmetric function are symmetric. For $i\ne1$ the influence of $i$ is the average of its influences in the two restrictions of $x_1$, so $\I[f]=\I_1[f]+\frac12(\I[f_{1\to1}]+\I[f_{1\to-1}])$. Induct on $n$, starting from constants.
\end{proof}

For symmetric $f$ on $n$ variables, the restrictions to $x_1=\pm1$ have coefficients $\wh f_s\pm\wh f_{s+1}$ on the $\binom Ns$ sets of size $s$ in the remaining coordinates, $N=n-1$. With $\varphi(u)=u\log_2(1/u)$ this gives the exact formulas
\begin{equation}\label{eq:delta}
 \Delta_1=\sum_{s=0}^{N}\binom Ns\,\mathcal D(\wh f_s,\wh f_{s+1}),\qquad \I_1=\sum_{s=0}^N\binom Ns\wh f_{s+1}^{\,2},
\end{equation}
\[
 \mathcal D(a,b)=\varphi(a^2)+\varphi(b^2)-\tfrac12\bigl[\varphi((a+b)^2)+\varphi((a-b)^2)\bigr].
\]
For $\mathrm{OR}_3$ one finds $\Delta_1/\I_1=12-\frac92\log_23$.

\begin{proof}[Proof of Theorem~\ref{thm:restriction}]
We work in natural units, $\mathcal D_e=(\ln2)\mathcal D$, and divide by $\ln2$ at the end.

\emph{Step 1: a tangent bound.} Let $h(t)=\mathcal D_e(\sqrt t,1)$ for $t\ge0$. Direct differentiation gives
\[
 h''(t)=\frac{1}{2t^{3/2}}\ln\Bigl|\frac{\sqrt t-1}{\sqrt t+1}\Bigr|<0,
\]
$h'$ is continuous at $t=1$ and tends to $0$ from above, so $h$ is increasing and concave and $(x,y)\mapsto\mathcal D_e(\sqrt x,\sqrt y)=y\,h(x/y)$ is concave on the quadrant. Every tangent gives $\mathcal D_e(a,b)\le h'(t)a^2+[h(t)-th'(t)]b^2$. Let $\mu_s=\binom Ns\wh f_{s+1}^{\,2}/\I_1$, a probability law on $\{0,\dots,N\}$. Taking the tangent at $t_s=(n-s)/s$ for the term $s\ge1$ of~\eqref{eq:delta}, and using $\binom N{s+1}/\binom Ns=(N-s)/(s+1)$ to move the $a^2$ part to the neighbouring index, one obtains
\[
 \frac{(\ln 2)\Delta_1}{\I_1}\le \frac{\mathcal D_e(\wh f_0,\wh f_1)}{\I_1}+\mu_0\kappa(\tfrac1n)+\sum_{s=1}^N\mu_s\Bigl[\Psi\bigl(\tfrac sn\bigr)+\kappa\bigl(\tfrac{s+1}n\bigr)-\kappa\bigl(\tfrac sn\bigr)\Bigr],
\]
where $\Psi(x)=h\bigl(\frac{1-x}x\bigr)$ and $\kappa(x)=\frac{1-x}xh'\bigl(\frac{1-x}x\bigr)$, extended continuously by $\kappa(0)=1$, $\kappa(1)=0$. Since $\kappa$ is uniformly continuous on $[0,1]$, the sum of the $\kappa$-differences tends to $0$ uniformly in $f$ as $n\to\infty$.

\emph{Step 2: three facts about $g=D_1f$.} (a) Let $m=|\wh f_1|=|\E g|$. Since $|g|=g^2$, $m\le\I_1$. By~\eqref{eq:alternating} and Lemma~\ref{lem:unimodal} applied to the binomial weights, $m\le b_N:=\max_j\binom Nj2^{-N}=O(N^{-1/2})$. From $h(t)\le3+\ln(1+t)$ and $|\wh f_0|\le1$ one gets $\mathcal D_e(\wh f_0,\wh f_1)/\I_1\le b_N[3+\ln(1+b_N^{-2})]\to0$, and $\mu_0=m^2/\I_1\le b_N\to0$. (b) Consecutive values of $g$ on adjacent weights have nonpositive product, so $\|D_jg\|_2^2\ge\frac12\|g\|_2^2$ for each $j$, and summing gives $\E_\mu s\ge N/2$. (c) By Theorem~\ref{thm:stab}, $\sum_s\mu_s\lambda^s\le\Phi\bigl(N(1-\lambda^2)/4\bigr)$ for $0\le\lambda\le1$.

\emph{Step 3: the limit law.} Take any sequence $n\to\infty$ with nonconstant symmetric $f_n$, and pass to a subsequence along which the law of $s/n$ under $\mu$ converges weakly to a law $\nu$ on $[0,1]$. By (b), $\E_\nu X\ge\frac12$. Taking $\lambda=e^{-t/n}$ in (c) and letting $n\to\infty$,
\begin{equation}\label{eq:laplace}
 \E_\nu e^{-tX}\le\Phi(t/2)\qquad(t>0),
\end{equation}
and in particular $\nu(\{0\})=0$. The function $\Psi$ is unbounded at $0$ but satisfies $\Psi(x)\le3+\ln(1/x)$, and (c) with $\lambda=e^{-1/m}$ gives $\mu[s\le m]\le e\,\Phi\bigl(N(1-e^{-2/m})/4\bigr)\le C\sqrt{m/N}$ for an absolute $C$. Summing over dyadic ranges, the contribution of $1\le s\le\delta n$ is $O(\sqrt\delta(1+\ln(1/\delta)))$ uniformly in $n$. This uniform integrability, Step~2(a) and weak convergence give $\sum_{s\ge1}\mu_s\Psi(s/n)\to\int\Psi\,d\nu$. Hence $\limsup(\ln2)\Delta_1/\I_1\le\sup_\nu\int\Psi\,d\nu$, the supremum over laws on $(0,1]$ with $\E_\nu X\ge\frac12$ and~\eqref{eq:laplace}.

\emph{Step 4: a Laplace majorant.} For $0<x\le1$,
\begin{equation}\label{eq:majorant}
 \Psi(x)\le U(x):=\Ei(40x)+\frac{553}{100}-\frac{273}{50}x+\frac{89}{50}e^{-20x}.
\end{equation}
This is verified by outward-rounded interval arithmetic on $11{,}900$ cells covering $[10^{-4},1]$, using that $\Psi$ and $U$ are decreasing, with minimum slack $0.0229$; for $0<x\le10^{-4}$ one uses $\Psi(x)\le3-\ln x$, $\Ei(40x)\ge-\gamma-\ln(40x)$ and $e^{-20x}\ge1-20x$. Since $\Ei(40x)=\int_{40}^\infty e^{-tx}\frac{dt}t$, Tonelli,~\eqref{eq:laplace} and $\E_\nu X\ge\frac12$ give
\[
 \int\Psi\,d\nu\le\int_{40}^\infty\Phi(t/2)\frac{dt}t+\frac{553}{100}-\frac{273}{100}+\frac{89}{50}\Phi(10).
\]
Finally $\Phi(t/2)=\frac2\pi\int_0^{\pi/2}e^{-t\sin^2u}du$. Using $\sin u\ge\frac{23}{24}u$ on $[0,\frac12]$ and $\sin u\ge\frac{23}{48}$ on $[\frac12,\frac\pi2]$,
\[
 \Phi(t/2)\le\frac{24}{23\sqrt{\pi t}}+e^{-t(23/48)^2},\qquad
 \int_{40}^\infty\Phi(t/2)\frac{dt}t\le\frac{48}{23\sqrt{40\pi}}+\frac{e^{-z}}z,
\]
with $z=40(23/48)^2$, by $\Ei(z)\le e^{-z}/z$. Collecting terms and dividing by $\ln2$ gives $C_*$, and a positive series for $I_0(10)$ certifies $\Phi(10)<0.127833337164$, so $C_*<4.637$.
\end{proof}

\section{Computations}\label{sec:computations}

The repository contains a script that recomputes every number stated above and prints a pass or fail line for each. Over all symmetric functions with $n\le16$ (exact integer Fourier coefficients, floating-point logarithms) it finds:
\begin{center}
\begin{tabular}{@{}ll@{}}
\toprule
quantity & maximum over symmetric $f$, $n\le16$\\
\midrule
$\HH/\I$ & $3.93027$, at AND/OR with $n=16$ \\
$L/\I$ & $2.51755$ \\
$B/\I$ & $1.80986$ \\
$\Delta_1/\I_1$ & $4.86767$, at $\mathrm{OR}_3$ \\
\bottomrule
\end{tabular}
\end{center}
In every dimension $2\le n\le16$ the maximum of $\HH/\I$ is attained by AND and OR. The script also checks the integral $C_N$ by adaptive quadrature for $N\le200$ and for $N$ up to $10^6$, Lemma~\ref{lem:bessel} on dense grids, Theorem~\ref{thm:stab} on random symmetric derivatives with $n=20,40,80$, the closed form of Theorem~\ref{thm:lower}, and the constants $C_{\rm sym}$ and $C_*$ at $40$ digits. A second script reproduces the interval certificate of~\eqref{eq:majorant}. These computations support but are not part of the proofs, except for the certificate of~\eqref{eq:majorant}.

\section{Questions}

\begin{question}
Is the optimal constant for symmetric functions equal to $4$?
\end{question}

\begin{question}\label{q:cutoff}
Find an explicit $n_0$ with $\Delta_1\le4.8723\,\I_1$ for every symmetric $f$ on $n\ge n_0$ variables. Together with a computation for $n<n_0$, Lemma~\ref{lem:induction} would then give $\HH\le4.8723\,\I$ for all symmetric $f$.
\end{question}

\begin{question}
Can the level bound $L\le\log_2(4+e)\I$ be improved for symmetric functions? Every gain there lowers $C_{\rm sym}$ directly.
\end{question}

\section*{Acknowledgements and funding}
The final manuscript was audited by Claude Opus 5.5 (Anthropic), which checked the proofs, recomputed the numerical constants in high-precision arithmetic, and reran the finite computations in the repository.

The research presented in this paper was supported by the European Research
Council (ERC) under the European Union's Horizon Europe research and
innovation programme (grant agreement No.~101041711), by the Simons
Foundation as part of the Collaboration on the Mathematical and Scientific
Foundations of Deep Learning, by Heights Labs, by Convex Nexus Capital, by
the Israel Science Foundation (grant number 2258/19), and by the Israel
Science Foundation (ISF Grant 4101/25). Computational resources were
provided in part by Google Cloud through its education credits program
(credit 530550103).

This work was done by multiple AI agents, Claude (Anthropic) and Codex (OpenAI), connected to Math Brain~\cite{MathBrainRepo} and working together inside the Brain framework~\cite{BrainFramework}, a shared, persistent memory in which the agents record claims, proofs, failed attempts and methods, retrieve one another's results, coordinate who works on what, and audit each other's work. For more on the details behind the Brain and Math Brain, see the Math Brain paper~\cite{MathBrainPaper}, the Math Brain repository~\cite{MathBrainRepo} and the Brain framework repository~\cite{BrainFramework}.

\section*{Declaration of competing interest}
The author holds equity in Heights Labs and in Convex Nexus Capital and
has a management role at Convex Nexus Capital. These entities are included
in the funding declaration above.

\begin{samepage}
\section*{Declaration of generative AI and AI-assisted technologies}
The paper comes out of Math Brain, an ongoing project on the use of AI systems for research mathematics, and its results were developed with the Brain and Math Brain. In November 2026, Claude Opus 5.5 (Anthropic) and a team of Codex agents running OpenAI's GPT-6.1, all connected to the Brain, retrieved the spectral-level entropy lemma and related methods from the author's earlier work through Math Brain, derived the split of the entropy into level and within-level parts and the noise-stability estimate, proved the Bessel-function inequalities and the asymptotic restriction estimate, carried out the finite computations, and searched the literature. Separate agents refereed the proof of the main theorem and of the asymptotic restriction estimate. Claude Opus 5.5 coordinated the agents, assembled the proofs, prepared the manuscript, and audited it, including the proofs, the numerical constants and the finite computations. These checks are not peer review. The author directed the work and is responsible for the content of the paper and for decisions concerning its dissemination.
\par
\end{samepage}

\section*{Data availability}
\url{https://github.com/yspennstate/fourier-entropy-influence-symmetric}

The paper is part of the Math Brain project, whose public repository is \url{https://github.com/yspennstate/math-brain-public}~\cite{MathBrainRepo}.

\begin{thebibliography}{99}
\bibitem{ACKSW21}
S.~Arunachalam, S.~Chakraborty, M.~Kouck\'y, N.~Saurabh and R.~de~Wolf,
\emph{Improved bounds on Fourier entropy and min-entropy},
ACM Trans. Comput. Theory \textbf{13} (2021), no.~4, Article~22, 40~pp.
\href{https://doi.org/10.1145/3470860}{doi:10.1145/3470860}.

\bibitem{BCV16}
J.-B.~Baillon, R.~Cominetti and J.~Vaisman,
\emph{A sharp uniform bound for the distribution of sums of Bernoulli trials},
Combin. Probab. Comput. \textbf{25} (2016), no.~3, 352--361.
\href{https://doi.org/10.1017/S0963548315000127}{doi:10.1017/S0963548315000127}.

\bibitem{CKLS16}
S.~Chakraborty, R.~Kulkarni, S.~V.~Lokam and N.~Saurabh,
\emph{Upper bounds on Fourier entropy},
Theoret. Comput. Sci. \textbf{654} (2016), 92--112.
\href{https://doi.org/10.1016/j.tcs.2016.05.006}{doi:10.1016/j.tcs.2016.05.006}.

\bibitem{FK96}
E.~Friedgut and G.~Kalai,
\emph{Every monotone graph property has a sharp threshold},
Proc. Amer. Math. Soc. \textbf{124} (1996), no.~10, 2993--3002.
\href{https://doi.org/10.1090/S0002-9939-96-03732-X}{doi:10.1090/S0002-9939-96-03732-X}.

\bibitem{KMS12}
N.~Keller, E.~Mossel and T.~Schlank,
\emph{A note on the entropy/influence conjecture},
Discrete Math. \textbf{312} (2012), no.~22, 3364--3372.
\href{https://doi.org/10.1016/j.disc.2012.07.031}{doi:10.1016/j.disc.2012.07.031}.

\bibitem{OD14}
R.~O'Donnell,
\emph{Analysis of Boolean Functions},
Cambridge University Press, Cambridge, 2014.

\bibitem{OWZ11}
R.~O'Donnell, J.~Wright and Y.~Zhou,
\emph{The Fourier entropy--influence conjecture for certain classes of Boolean functions},
in: ICALP 2011, Lecture Notes in Comput. Sci. \textbf{6755}, Springer, Heidelberg, 2011, 330--341.
\href{https://doi.org/10.1007/978-3-642-22006-7_28}{doi:10.1007/978-3-642-22006-7\_28}.

\bibitem{Shm26}
Y.~Shmalo,
\emph{Fourier entropy of threshold functions and their compositions},
preprint, 2026.
\href{https://arxiv.org/abs/2610.16069}{arXiv:2610.16069}.

\bibitem{MathBrainRepo} Y.~Shmalo, \emph{Math Brain: public snapshot, code and tests}, 2026, \url{https://github.com/yspennstate/math-brain-public}.

\bibitem{MathBrainPaper} Y.~Shmalo, \emph{Math Brain: shared, auditable memory for multi-agent mathematical research}, October 2026, \url{https://github.com/yspennstate/math-brain-public/tree/main/papers/math-brain-2026-10-24}.

\bibitem{BrainFramework} Y.~Shmalo, \emph{Brain Framework}, public software repository, 2026, \url{https://github.com/yspennstate/brain-framework-public}.
\end{thebibliography}
\end{document}
